\section*{Chapter \ref{ch:g6:while-loops} --- Repeating with \texttt{while}}

\begin{solution}{exo:g6:while-loops:1}
5, 3 and 1. After 1, \code{n} becomes $-1$, and the \omterm{def:g4:algorithms-repeat-and-decide:condition}{condition} \code{n > 0} is
false.
\end{solution}

\begin{solution}{exo:g6:while-loops:2}
5 times: \code{n} becomes 3, 9, 27, 81 and 243. After 81 the test $81 < 100$
is still true, so the block runs once more; 243 is the last number printed.
\end{solution}

\begin{solution}{exo:g6:while-loops:3}
\begin{center}
\begin{tabular}{|c|c|c|}
\hline
turn & \code{total} & \code{i} \\ \hline
before & 0 & 1 \\
1 & 1 & 2 \\
2 & 5 & 3 \\
3 & 14 & 4 \\
4 & 30 & 5 \\ \hline
\end{tabular}
\end{center}
With \code{i} equal to 5 the test \code{i <= 4} is false: the \omterm{def:g4:programs-and-operating-system:program}{program} prints
30, which is $1 + 4 + 9 + 16$.
\end{solution}

\begin{solution}{exo:g6:while-loops:4}
\omcode{code/grade-6/05-while-loops/exercises/sum100.py}{1}{6}{The sum of 1 to
100.}\label{lst:g6:while-loops:sol-sum100}

It prints 5050.
\end{solution}

\begin{solution}{exo:g6:while-loops:5}
Nothing in the block changes \code{i}: it stays 0, and \code{i < 10} stays true.
Add \code{i = i + 1} at the end of the block, indented like the \code{print}.
\end{solution}

\begin{solution}{exo:g6:while-loops:6}
The 10 is missing: when \code{i} reaches 10, the test \code{i < 10} is false,
so 10 is never added. With \code{while i <= 10:} the \omterm{def:g4:programs-and-operating-system:program}{program} prints 55.
\end{solution}

\begin{solution}{exo:g6:while-loops:7}
\omcode{code/grade-6/05-while-loops/exercises/doubling.py}{1}{6}{Doubling past
1,000.}\label{lst:g6:while-loops:sol-doubling}

It prints \code{1024 after 10 doublings}: 512 is still not above 1,000, and
the next doubling gives 1,024.
\end{solution}

\begin{solution}{exo:g6:while-loops:8}
\omcode{code/grade-6/05-while-loops/exercises/digit_sum.py}{1}{6}{The sum of the
digits.}\label{lst:g6:while-loops:sol-digit-sum}
\end{solution}

\begin{solution}{exo:g6:while-loops:9}
\code{n} takes the \omterm{def:g6:variables-and-input:value}{values} 50, 25, 12, 6, 3 and 1: 6 turns. $25 // 2$ is 12 and
$3 // 2$ is 1, since \code{//} drops the remainder.
\end{solution}

\begin{solution}{exo:g6:while-loops:10}
The question is asked once before the \omterm{def:g4:algorithms-repeat-and-decide:loop}{loop} and again at the end of each turn.

\omcode{code/grade-6/05-while-loops/exercises/ask_yes.py}{1}{5}{Asking until the
answer is yes.}\label{lst:g6:while-loops:sol-ask}
\end{solution}

\begin{solution}{exo:g6:while-loops:11}
With $c$ coins a week, $c$ at least 1, \code{coins} grows by at least 1 at each
turn: after 50 turns at most, it is at least 50, and the test \code{coins < 50}
is false. With 0 coins a week, \code{coins} stays 0 for ever: an \omterm{def:g6:while-loops:infinite}{infinite loop}.
\end{solution}

\begin{solution}{exo:g6:while-loops:12}
The \omterm{def:g4:algorithms-repeat-and-decide:condition}{condition} gains a second part with \code{and}: the \omterm{def:g4:algorithms-repeat-and-decide:loop}{loop} also stops after
the seventh guess. After the \omterm{def:g4:algorithms-repeat-and-decide:loop}{loop}, an \code{if} tells why it stopped.

\omcode{code/grade-6/05-while-loops/exercises/guess_limit.py}{1}{15}{At most 7
guesses.}\label{lst:g6:while-loops:sol-guess}
\end{solution}

\begin{solution}{pb:g6:while-loops:1}
\textbf{1.} 6, 3, 10, 5, 16, 8, 4, 2, 1: 8 steps; the highest number is 16.

\textbf{2.} 7, 22, 11, 34, 17, 52, 26, 13, 40, 20, 10, 5, 16, 8, 4, 2, 1: 16
steps; the highest is 52.

\textbf{3.} A doubling of 1 is even, its half is again a doubling of 1, and so
on down to 1: the game only halves. 64, 32, 16, 8, 4, 2, 1: 6 steps.

\textbf{4.} \code{steps} is a \omterm{def:g6:while-loops:counter}{counter} (one more at each turn). \code{highest}
keeps the largest number seen so far; it starts at the starting number because
the start itself may be the highest, as for 16.

\textbf{5.} It prints \code{steps: 8 highest: 16} for 6 and \code{steps: 16
highest: 52} for 7, as found by hand.

\textbf{6.} 27, a small start, takes 111 steps and climbs as high as 9,232
before coming down:

\omoutput{code/grade-6/05-while-loops/outputs/collatz.txt}{The start
27.}\label{out:g6:while-loops:collatz}

\textbf{7.} The \code{while} \omterm{def:g4:algorithms-repeat-and-decide:loop}{loop} of \cref{lst:g6:while-loops:collatz} goes
inside a \code{for} \omterm{def:g4:algorithms-repeat-and-decide:loop}{loop} over the starts, with a record of the best start:

\omcode{code/grade-6/05-while-loops/listings/collatz_record.py}{1}{16}{The start
below 100 with the most steps.}\label{lst:g6:while-loops:collatz-record}

\omoutput{code/grade-6/05-while-loops/outputs/collatz_record.txt}{97 takes 118
steps.}\label{out:g6:while-loops:collatz-record}

\textbf{8.} Running \cref{lst:g6:while-loops:collatz} with 97: the highest
number reached is 9,232, the same as for 27: the two paths meet at 94 and go
on together.

\textbf{9.} No. It proves that the starts from 1 to 99 reach 1, and nothing
about 100 or any larger start.

\textbf{10.} Its \code{while} \omterm{def:g4:algorithms-repeat-and-decide:loop}{loop} would never end for that start: an \omterm{def:g6:while-loops:infinite}{infinite
loop}, and the \omterm{def:g4:programs-and-operating-system:program}{program} would print nothing.

\textbf{11.} ``Checked'' means that each number tried reached 1. ``Proved''
would mean an argument showing that every number, without exception, reaches
1. There are infinitely many whole numbers, so no \omterm{def:g1:what-is-a-computer:computer}{computer} can try them all;
only a proof could settle the question, and nobody has found one.

\textbf{12.} \textbf{Result.} Among the starts below 100, 97 takes the most
steps to reach 1: 118 steps.
\end{solution}
